\documentclass[12pt, a4paper]{article}

\usepackage[margin=1in]{geometry}
\usepackage{hyperref}
\usepackage{titlesec}
\usepackage{enumitem}
\usepackage{booktabs}
\usepackage{xcolor}
\usepackage{microtype}
\usepackage{amsmath}
\usepackage{fancyhdr}

\hypersetup{
    colorlinks=true,
    linkcolor=blue!70!black,
    urlcolor=blue!70!black,
}

\pagestyle{fancy}
\fancyhf{} % Clears all default header and footer fields
\fancyhead[R]{sujeet.banerjee@gmail.com
\\
sujeet.banerjee.professional@gmail.com} % Sets the right header
\fancyhead[L]{Agentic AI Bootcamp}
\fancyfoot[R]{\thepage} % Keeps the page number at the bottom center
\fancyfoot[L]{Advanced: PDF Data Extraction}
\renewcommand{\headrulewidth}{0.4pt} % Adds a decorative line under the header (optional)
\renewcommand{\footrulewidth}{0.4pt}

\title{\textbf{Advanced Strategies for PDF Text Extraction and Data Cleaning in RAG Systems}}
\author{Detailed Notes - Agentic AI Bootcamp}
\date{}


\begin{document}

\maketitle
\tableofcontents
\newpage

\section{Introduction}
In Retrieval-Augmented Generation (RAG) pipelines, the quality of the generated response is strictly upper-bounded by the quality of the retrieved context. PDFs remain the de facto standard for enterprise documents (reports, research papers, manuals), but their format is fundamentally presentation-oriented, not semantically structured. This document outlines advanced strategies for extracting, parsing, and cleaning text from various PDF types before the chunking stage.

\section{Categorization of PDF Documents}
Before applying extraction techniques, it is essential to classify the PDF into one of two primary categories:
\begin{itemize}
    \item \textbf{Text-Based (Native) PDFs:} Generated directly from digital sources (e.g., MS Word, LaTeX). They contain character encodings and explicit text layers.
    \item \textbf{Scanned (Image-Based) PDFs:} Created via physical scanners. They contain no text layer, only rasterized images of pages.
\end{itemize}

\section{Extraction Strategies for Text-Based PDFs}
Native PDFs require libraries that can read the internal text streams.

\subsection{Standard Text Extraction}
\begin{itemize}
    \item \textbf{PyMuPDF (fitz):} Highly performant, offers bounding box access, and preserves reading order reasonably well. Ideal for extracting text, images, and basic layout information.
    \item \textbf{pdfplumber:} Excellent for detailed character-level bounding box analysis. Slower than PyMuPDF but superior for debugging complex layouts and extracting tabular data.
\end{itemize}

\subsection{Layout-Aware Parsing}
Standard extraction often fails on multi-column layouts, interleaving text across columns.
\begin{itemize}
    \item \textbf{Heuristic Methods:} Sorting text blocks by their $(x, y)$ coordinates.
    \item \textbf{Machine Learning Models:} Utilizing layout analysis models (e.g., LayoutLMv3, YOLOX-based parsers like Unstructured) to classify blocks into \texttt{Title}, \texttt{Text}, \texttt{List}, \texttt{Table}, and \texttt{Figure}.
\end{itemize}

\section{Extraction Strategies for Scanned PDFs (OCR)}
For scanned documents, Optical Character Recognition (OCR) is mandatory.

\subsection{Pre-processing Images}
OCR accuracy drops significantly with poor image quality. Pre-processing steps include:
\begin{itemize}
    \item \textbf{Binarization:} Converting grayscale/color images to black and white (e.g., Otsu's thresholding, adaptive thresholding).
    \item \textbf{Deskewing:} Rotating the image to correct scanning tilt.
    \item \textbf{Denoising:} Applying Gaussian blurs or median filters to remove background artifacts.
\end{itemize}

\subsection{OCR Engines}
\begin{itemize}
    \item \textbf{Tesseract OCR:} Open-source, supports many languages. Requires rigorous image pre-processing for optimal results.
    \item \textbf{EasyOCR / PaddleOCR:} Deep learning-based, handles varied fonts, orientations, and complex backgrounds better than Tesseract.
    \item \textbf{Cloud APIs:} AWS Textract, Google Cloud Document AI, and Azure Document Intelligence provide state-of-the-art accuracy, including built-in layout and table understanding.
\end{itemize}

\section{Specialized Extraction Techniques}

\subsection{Table Extraction}
Tables contain dense relational data that standard text extraction flattens into unreadable strings.
\begin{itemize}
    \item \textbf{Tools:} Camelot, Tabula, or OCR-based table parsers (Table Transformer).
    \item \textbf{Representation:} Convert extracted tables into Markdown, HTML, or CSV string formats before chunking to preserve row/column relationships for the LLM.
\end{itemize}

\subsection{Optical Mark Recognition (OMR)}
Used for forms with checkboxes or radio buttons.
\begin{itemize}
    \item \textbf{Strategy:} Detect contour boxes using OpenCV. Calculate the pixel density (ratio of black to white pixels) inside the box. A density above a specific threshold implies a ``checked'' state.
    \item \textbf{Mapping:} Map the checked state to adjacent text to yield key-value pairs (e.g., \texttt{Is\_Employed: True}).
\end{itemize}

\subsection{Figures and Charts}
\begin{itemize}
    \item \textbf{Extraction:} Use PyMuPDF to extract images.
    \item \textbf{Vision-Language Models (VLMs):} Pass extracted images through models like LLaVA or GPT-4o to generate descriptive captions, injecting these captions back into the text stream.
\end{itemize}

\section{Data Cleaning Post-Extraction}
Raw extracted text is rarely ready for embedding and chunking. A robust cleaning pipeline is critical.

\subsection{Artifact Removal}
\begin{itemize}
    \item \textbf{Headers/Footers \& Page Numbers:} Often interrupt sentence flow across page boundaries. Remove using regex (e.g., \verb|^\d+$|) or by ignoring text within the top/bottom 10\% of the page bounding box.
    \item \textbf{Watermarks:} Filter out repetitive strings that have low opacity or span diagonally across bounding boxes.
\end{itemize}

\subsection{Text Normalization}
\begin{itemize}
    \item \textbf{De-hyphenation:} Words split across lines (e.g., ``docu-\textbackslash nment'') must be rejoined. Use regex: \verb|re.sub(r'-\n\s*', '', text)| (with dictionary checks to avoid merging legitimate hyphenated words).
    \item \textbf{Whitespace Normalization:} Collapse multiple spaces, tabs, and spurious newlines into single spaces.
    \item \textbf{Unicode Normalization:} Convert ligatures (e.g., ``\mbox{fi}'', ``\mbox{fl}'') to standard characters using NFKC normalization.
\end{itemize}

\subsection{Formatting and Metadata Enrichment}
\begin{itemize}
    \item \textbf{Markdown Conversion:} Where possible, convert headings into Markdown (e.g., \texttt{\# Heading}) to provide structural cues for chunking algorithms (like LangChain's \texttt{MarkdownHeaderTextSplitter}).
    \item \textbf{Metadata Tagging:} Append source metadata (page number, document ID, section title) to the cleaned text blocks. This is vital for downstream Graph-RAG or precise citation generation.
\end{itemize}

\section{Conclusion}
Thorough PDF extraction and data cleaning are non-negotiable precursors to effective RAG. By treating native and scanned PDFs distinctly, deploying layout-aware parsing, utilizing OMR for forms, and rigorously normalizing the extracted text, we ensure that the subsequent chunking phase operates on coherent, high-fidelity semantic units.

\end{document}
